\section{Introduction}
\label{sec:introduction}

The squared geodesic distance is smooth before the cut locus, but optimal
transport and its supporting functions naturally reach that boundary.
The weak Ma--Trudinger--Wang condition controls the cost Hessian only in
directions perpendicular to a velocity segment. We prove that this
transverse condition already forces global convexity and supporting geometry
at arbitrary minimizing endpoints, including conjugate cut points.

\subsection{Weak MTW and convexity}

Let $(M,g)$ be a smooth, connected, compact Riemannian manifold without
boundary, of dimension $n\geq2$. Write $d$ for its geodesic distance,
$\exp_x:T_xM\to M$ for its exponential map, and $|\cdot|_x$ and
$\langle\cdot,\cdot\rangle_x$ for the norm and inner product of $g_x$.
The open tangent injectivity domain is
\begin{equation}
 I(x)=\left\{v\in T_xM:\ \text{there is }a>1\text{ such that }
 d(x,\exp_x(av))=a|v|_x\right\}.
 \label{eq:injectivity-domain}
\end{equation}
Its geodesics remain minimizing beyond time one. Put $c(x,y)=d(x,y)^2/2$.
Our MTW convention is
\begin{equation}
 \mathfrak S_{(x,v)}(\xi,\eta)
 =-\frac32\left.
 \frac{\partial^4}{\partial s^2\partial t^2}
 c\bigl(\exp_x(t\xi),\exp_x(v+s\eta)\bigr)
 \right|_{s=t=0},\qquad v\in I(x).
 \label{eq:mtw-definition}
\end{equation}
The cost is smooth near the pairs used here, as recalled in
\cref{sec:preliminaries}. Throughout the paper we assume precisely
\begin{equation}
 \langle\xi,\eta\rangle_x=0
 \quad\Longrightarrow\quad
 \mathfrak S_{(x,v)}(\xi,\eta)\geq0
 \qquad(x\in M,\ v\in I(x),\ \xi,\eta\in T_xM).
 \label{eq:weak-mtw}
\end{equation}
No condition on nonorthogonal pairs, strict lower bound, nonfocality,
or prior convexity of $I(x)$ is assumed.
Write $G_x=\{p\in T_xM:d(x,\exp_xp)=|p|_x\}$ for the closed
minimizing domain.

\begin{theorem}[Convexity of injectivity domains]\label{thm:main}
If $(M,g)$ satisfies \eqref{eq:weak-mtw}, then $I(x)$ is convex for
every $x\in M$. More precisely, for every $v_0,v_1\in I(x)$ and
$\theta\in[0,1]$,
\[
 (1-\theta)v_0+\theta v_1\in I(x).
\]
\end{theorem}

In fact the closed minimizing domain $G_x$ is convex as well
(\cref{geo:convexity}). The conclusion is ordinary linear convexity;
no strict convexity or regularity of the boundary is asserted.

Villani proposed that weak MTW should force convexity of every tangent
injectivity domain \cite[Section~4.2]{Villani2011}. In the compact setting
of \cref{thm:main}, the conjecture imposes neither nonfocality nor prior
convexity. Figalli, Gallou\"et and Rifford proved the implication for
nonfocal manifolds \cite[Theorem~1.7]{FGR2014}, where minimizing geodesics
reach cut strictly before their first conjugate point. Here and below the
numbering in that citation refers to the published article.
\Cref{thm:main} removes this restriction: conjugate cut endpoints are
allowed in every dimension $n\geq2$.

Loeper and Villani connected MTW to cut geometry, obtaining convexity
under strong MTW and nonfocality and treating some focal configurations
under stronger geometric conditions \cite{LoeperVillani2010}; see also
\cite[Section~3.5.2]{Figalli2010}. Their work and the nonfocal theorem of
\cite{FGR2014} distinguish the two obstructions in the present problem.
The cost may cease to be smooth along a proposed velocity segment, and
weak MTW controls only directions perpendicular to that segment.
Neither difficulty can be removed by applying smooth-cost concavity
beyond its domain of definition.

\subsection{Global support and intermediate geometry}

A \emph{potential} is a function
$u(x)=\max_y\{-c(x,y)-v(y)\}$ for a continuous datum $v:M\to\mathbb R$;
replacing $v$ by the dual transform
$v(y)=\max_x\{-c(x,y)-u(x)\}$ gives a dual pair with nonnegative gap
$u(x)+c(x,y)+v(y)$. Such functions are Lipschitz and semiconvex.
We use their ordinary semiconvex subdifferentials $\partial u(x)$,
identifying vectors and covectors by the metric. Define
\[
 \Gamma_u=\{(x,p):p\in\partial u(x)\},\qquad
 Q_tf(z)=\min_x\left\{f(x)+\frac{c(x,z)}t\right\}.
\]

\begin{theorem}[Global support and intermediate geometry]
\label{thm:foundation}
Under \eqref{eq:weak-mtw}, every potential $u$, with its dual potential $v$,
has the following properties.
\begin{enumerate}[label=\textup{(\roman*)}]
\item Every $p\in\partial u(x)$ belongs to $G_x$ and satisfies
\[
 u(x')\geq u(x)+c(x,\exp_xp)-c(x',\exp_xp)
 \qquad(x'\in M).
\]
For every $0<t<1$, the map $F_t:\Gamma_u\to M$,
$F_t(x,p)=\exp_x(tp)$, is a homeomorphism; $tp\in I(x)$ and $x$
is the unique global minimizing pole of $Q_tu$ at $F_t(x,p)$.
\item For every dual pair $(u,v)$, every $x$ and every $r$, the lifted section
\[
 \mathcal S_x(r)=\left\{p\in G_x:
 u(x)+\frac{|p|^2}{2}+v(\exp_xp)\leq r\right\}
\]
is compact and convex. If it is nonempty and $r\geq0$, then
\[
 \operatorname{diam}\mathcal S_x(r)^2
 \leq8\left(r+\osc_{p\in\mathcal S_x(r)}v(\exp_xp)\right).
\]
\item For $0<t<1$, one has $Q_tu=-Q_{1-t}v\in C^{1,1}(M)$.
Its $C^{1,1}$ seminorm and the Lipschitz constants of the pole and momentum
maps $F_t^{-1}$ depend only on $M,g,t$. There are also $r_t,b_t>0$,
depending only on $M,g,t$, such that, for $z=F_t(x,p)$,
\[
 u(x')+\frac{c(x',z)}t
 \geq u(x)+\frac{c(x,z)}t+b_t d(x,x')^2
 \quad\text{if }d(x,x')<r_t.
\]
These constants can be chosen uniformly on compact subintervals of $(0,1)$.
\end{enumerate}
No density hypothesis is required.
\end{theorem}

The detailed statements and proofs are \cref{geo:support,geo:sections,geo:intermediate}.
The distinction between intermediate times and time one is essential:
only the support and minimizing-velocity assertions pass to time one.
For instance, $u(x)=-c(x,y_0)$ has many poles with the same final endpoint.

\subsection{Supporting geometry and earlier methods}

Ma, Trudinger and Wang introduced a strictly positive curvature condition
in optimal-transport regularity theory \cite{MTW2005}; Trudinger and Wang
treated its degenerate, nonnegative form, usually denoted A3w
\cite[Section~1]{TW2009}. Loeper established necessity for continuity in
the smooth-cost setting and developed the supporting-mountain principle
\cite{Loeper2009}. Kim and McCann gave its geometric formulation and a
double-mountain principle under smooth-cost and illuminated-set
hypotheses \cite[Theorem~4.10]{KM2010}. For squared distance on a compact
manifold, Figalli, Rifford and Villani obtained the global supporting
principle, including endpoint closure, from weak MTW together with convex
injectivity domains \cite[Lemma~3.1]{FRV2011}. The issue here is to derive
that convexity rather than assume it.

Several ingredients precede the MTW theory. The shortened-distance
support inequality is \cite[Claim~2.4]{CMS2001}. The endpoint-moving
Jacobi trial used at a possible conjugate cut is related to the
second-variation argument in \cite[Proposition~2.5]{CMS2001}; the
transverse MTW reduction and the joint radial--tangential growth
contradiction are proved here. Intermediate injectivity at
differentiability points and inverse estimates on compact subsets of a
positive-Hessian set appear in \cite[Lemma~5.3 and Claim~5.6]{CMS2001}.
Our conclusion concerns the entire ordinary subgradient graph and gives
bounds depending only on the geometry and the intermediate time.

The graph coordinates in the optional degree argument and the infimum
regularization also have classical
antecedents. Moreau developed quadratic inf-convolution and proximal
decomposition \cite[Sections~3--5, 7 and~12]{Moreau1965}; Minty's
monotone-graph coordinates are presented in
\cite[Proposition~1.1]{AA1999} and adapted to transport-plan
rectifiability in \cite[Section~2]{MPW2012}. We prove the local
semiconvex graph representation needed here. These constructions do not
by themselves supply continuation through cut or the global
support theorem. The resulting density-free geometry is distinct from
the additional measure estimates needed for regularity of transport
maps.

\subsection{The continuation argument}

\Cref{sec:preliminaries} develops the minimizing-branch Hessian and index
form without nonfocality. Proper radial shortening supplies smooth
distance supports, and the divided Hessian decreases strictly with
duration. These estimates also control the uniform constants in the
intermediate conclusions.

The proof works on the entire ordinary subgradient graph of an arbitrary
potential. A target is active at $x$ if it attains the maximum defining
$u(x)$; an active log is any minimizing velocity to that target.
Every ordinary subgradient is a convex combination of finitely many
active logs. All minimizing logarithms of each active target are retained.
These are finite local representations; the
potential itself need not have finitely many targets.

At a first failing scale, limits from earlier scales show that the
relevant finite convex hulls are still minimizing. The argument of
\cref{sec:conjugacy} excludes conjugacy there. Transverse concavity
confines a hypothetical conjugacy kernel to the segment direction;
the radial Hessian identity then contradicts the growth forced by a
Jacobi trial field.

To exclude local collisions of the graph projection, the finite active
representations are compared before taking limits. The first-order inequalities force the positive-weight limiting velocities
into an affine hyperplane perpendicular to the collision direction.
The linear terms are then canceled exactly at the finite stage,
before division by the squared collision distance and passage to the limit.
Weak MTW and strict radial comparison give the contradiction.
Invariance of domain now makes the projection a local homeomorphism.
Compactness and a homotopy to the short-time graph homeomorphism make
this covering one-sheeted. Capturing every minimizing logarithm
excludes ordinary cuts and continues the construction to every time
strictly below one. Closure yields global support at time one; convex
lifted sections and two-sided intermediate supports give the remaining
conclusions of \cref{thm:foundation}.

An optional alternative in \cref{sec:degree-alternative} replaces local
injectivity by local quadratic growth and openness. A Jacobian-sign
calculation, degree one, and the area formula then give the same
conditional globalization. It reuses the common setup and is not
needed for the main proof.
